% !TEX root =  master.tex
%%
%% @stud: sprachspezifische Anpassung
%%
\chapter*{Kurzfassung (Abstract)}
\addcontentsline{toc}{chapter}{Kurzfassung (Abstract)}

Diese Projektarbeit gibt einen strukturierten Überblick über rekurrente neuronale Netze und deren Einsatz bei der Verarbeitung sequenzieller Daten. Ausgehend von den mathematischen Grundlagen und zentralen Konzepten des maschinellen Lernens wird zunächst das Grundmodell einfacher rekurrenter Netze beschrieben. Darauf aufbauend wird das Training mittels \ac{BPTT} erläutert und in den Kontext der allgemeinen Gradientenberechnung in neuronalen Netzen eingeordnet.

Ein Schwerpunkt der Arbeit liegt auf den typischen Schwierigkeiten beim Training rekurrenter Architekturen. Insbesondere werden die Ursachen und Auswirkungen verschwindender sowie explodierender Gradienten analysiert. Als zentrale Gegenmaßnahmen werden Gradienten-Clipping sowie die erweiterten Architekturen \ac{LSTM} und \ac{GRU} vorgestellt und hinsichtlich ihrer Funktionsweise eingeordnet.

Die Arbeit verfolgt das Ziel, die theoretischen Grundlagen, die praktischen Herausforderungen und die wesentlichen Lösungsansätze im Umfeld rekurrenter neuronaler Netze kompakt und nachvollziehbar darzustellen. Damit schafft sie eine Grundlage für die Bewertung des Einsatzes dieser Modelle in Anwendungsfeldern wie Zeitreihenanalyse und Sprachverarbeitung.

